\section{給付}

\subsection{給付の種類・裁定・受給要件}

\begin{InyouJobunbox}
  \begin{itemize}
    \item 法第29条〜第31条
    \item 則第33条、第36条
  \end{itemize}
\end{InyouJobunbox}

\begin{fillblank}[年金1 2023 問題1(1)アイ]
  \noindent \textbf{◯DB法}
  \begin{lawstructure}
    \begin{lawarticle}{第三十条}{裁定}
      給付を受ける権利（以下「受給権」という。）は、その権利を有する者（以下「受給権者」という。）の請求に基づいて、（　　A　　）等が裁定する。
      
      \lawparagraph{2}
      （　　A　　）は、前項の規定により裁定をしたときは、遅滞なく、その内容を資産管理運用機関に通知しなければならない。

      \lawparagraph{3}
      資産管理運用機関又は（　　B　　）（以下「資産管理運用機関等」という。）は、第一項の規定による裁定に基づき、その請求をした者に給付の支給を行う。
      
    \end{lawarticle}
  \end{lawstructure}
\end{fillblank}

\begin{answerbox}
  \begin{enumerate}[A:]
    \item 事業主
    \item 基金
  \end{enumerate}
\end{answerbox}

\begin{explanationbox}
  「事業主等」とは、事業主と基金を指す。（法第29条）
\end{explanationbox}

\newpage

\subsection{給付の額１（総論）}

\begin{InyouJobunbox}
  \begin{itemize}
    \item 法第32条
    \item 解釈第3-1\textcircled{1}\textcircled{2}\textcircled{3}、解釈第4-1(3)
  \end{itemize}
\end{InyouJobunbox}

\begin{shortanswer}[年金1 2022 問題2(1)ア]
  給付の額に関する事項として、確定給付企業年金法第32 条において、「加入者期間又は当該
加入者期間における給与の額その他これに類するものに照らし、適正かつ合理的なものとし
て政令で定める方法により算定されたものでなければならず、かつ、特定の者について不当
に差別的なものであってはならない」と定められている。その取扱いについて、制度の目的
が老後の安定的所得の保障であることに鑑み、制度の目的を逸脱するものであってはならな
いこととして『確定給付企業年金法並びにこれに基づく政令及び省令について（法令解釈）』
に例示されているものを3 つ簡潔に入力しなさい。
\end{shortanswer}

\begin{answerbox}
  \begin{itemize}
    \item 資格喪失事由や資格喪失時の年齢等により給付の額に格差を設ける場合であって、給付の額の格差が過大であること
    \item 上記の場合で、早期に脱退した者の給付の額の方が有利であること
    \item 給付の額の算定方法に上限の設定が含まれている場合においても、早期に脱退した者の給付の額の方が有利であるなど
  \end{itemize}
\end{answerbox}



\begin{fillblank}[年金1 2025 問題1(1)]
  \noindent \textbf{◯法令解釈}
  \begin{lawstructure}
    \begin{lawarticle}{}{}
      \noindent \textbf{第３ 給付の額に関する事項}

      \lawparagraph{1}
      給付の額は、法第３２条第２項において「加入者期間又は当該加入者期間における
      給与の額その他これに類するものに照らし、適正かつ合理的なもの」により算定され
      たものでなければならず、かつ、「特定の者について不当に差別的なものであっては
      ならない」と規定されているところであるが、その取扱いは次のとおりとすること。

      \lawitem{\textcircled{1}}
給付の額は、（　　A　　）に応じて算定されるものであり、原則として、（　　A　　）
が長くなるにもかかわらず給付の額が減少するものであってはならないこと。こ
のため、障害給付金であっても、若年者に支給する額は年長者に支給する額に比し
て過大なものとならないこと。


      \lawitem{\textcircled{2}}
加入者間で給付の額に差を設ける場合にあっては、「短時間・有期雇用労働者及
び派遣労働者に対する（　　B　　）の禁止等に関する指針」を踏まえ、労働協約等
において、特定の職種に属する従業員や特定の学歴の従業員に係る給与及び退職金
等の労働条件が他の職種に属する従業員や他の学歴の従業員の労働条件とは別に規
定されているなど、給付の額に差を設けることにつき合理的な理由があること。

      \lawitem{\textcircled{3}}
制度の目的が（　　C　　）であることに鑑み、資格喪失事由や資格喪
失時の年齢等により給付の額に格差を設ける場合においても、給付の額の格差が過
大であること、（　　D　　）の給付の額の方が有利であることなど、制度の目
的を逸脱するものであってはならないこと。また、給付の額の算定方法に規則第２
５条第２号の方法に基づく上限の設定が含まれている場合においても、（　　D　　）
の給付の額の方が有利であることなど、制度の目的を逸脱するものであって
はならないこと。

      （以下略）

    \end{lawarticle}
  \end{lawstructure}  
\end{fillblank}

\begin{answerbox}
  \begin{enumerate}
    \item 加入者期間
    \item 不合理な待遇
    \item 老後の安定的所得の保証
    \item 早期に脱退した者
  \end{enumerate}
\end{answerbox}

\begin{truefalse}[年金1 2020 問題1(3)\textcircled{2}]
    \textbf{下線部の内容の正誤を判定し、誤っている場合は正しい内容に改めなさい。}
  \\

  加入者が負担する掛金のある制度において、加入者がその選択により掛金を負担しない場合
の給付の額は、掛金を負担する加入者に比べて、当該\underline{掛金の拠出額の下限予定利率に基づく元
利合計額に相当}する程度の差を設けるものであること。
\end{truefalse}

\begin{answerbox}
  \textbf{誤}

  掛金の拠出額に相当
\end{answerbox}

\begin{fillblank}[年金1 2019 問題1(1)A]
  \noindent \textbf{◯DB法}
  \begin{lawstructure}
    \begin{lawarticle}{第三十二条}{確定給付企業年金の実施}
      給付の額は、政令で定める基準に従い規約で定めるところにより算定した額とする。

      \lawparagraph{2}
      前項に規定する給付の額は、加入者期間又は当該加入者期間における給与の額その他これに類するものに照らし、適正かつ合理的なものとして政令で定める方法により算定されたものでなければならず、かつ、（　　A　　）であってはならない。
    \end{lawarticle}
  \end{lawstructure}  
\end{fillblank}

\begin{answerbox}
  \begin{enumerate}[A:]
    \item 特定の者について不当に差別的なもの
  \end{enumerate}
\end{answerbox}

\newpage

\subsection{給付の額２（給付の額の基準）}


\begin{InyouJobunbox}
  \begin{itemize}
    \item 令第23条
    \item 則第24条〜第24条の4
    \item 解釈第3-1\textcircled{8}\textcircled{9}、解釈第3-2、解釈第3-7、解釈第3-8
  \end{itemize}
\end{InyouJobunbox}

\begin{shortanswer}[年金1 2019 問題2(3)\textcircled{1}]
  一時金として支給する老齢給付金の額は、当該老齢給付金の全部を年金として支給するとし
た場合の老齢給付金のうち、保証期間について支給する給付現価相当額を上回らないものとし
なければならない。受託保証型確定給付企業年金でない確定給付企業年金において、この現価
相当額を計算するときの基礎となる予定利率および予定死亡率について簡記しなさい。
\end{shortanswer}

\begin{answerbox}
  ＜予定利率＞

次に掲げる率のうち最も低い率。
\begin{enumerate}
  \item 前回の財政計算の計算基準日以降の日における下限予定利率のうち、最も低い下限予定利率
  \item 老齢給付金支給開始要件を満たした時における1の率
  \item 加入者の資格を喪失したときにおける1の率
\end{enumerate}

＜予定死亡率＞

前回の財政計算で使用した率。
\end{answerbox}

\begin{shortanswer}[自作 2026/10/11]
  前問の現価相当額を計算するときの予定利率について、受託保証型確定給付企業年金である場合の定め方を簡記しなさい。
\end{shortanswer}

\begin{answerbox}
  契約者価額の計算に用いる予定利率
\end{answerbox}

\begin{truefalse}[年金1 H28 問題1(5)\textcircled{4}]
    \textbf{下線部の内容の正誤を判定し、誤っている場合は正しい内容に改めなさい。}
  \\  

  確定給付企業年金における遺族給付金の額は、 確定給付企業年金法施行令第 23 条第 3 項に該
当する場合を除き、老齢給付金の受給権者となった者が受給権の取得と同時に死亡した場合に
おいてその者の遺族に支給する遺族給付金の現価相当額が、当該老齢給付金の全部を年金とし
て支給するとした場合の\underline{老齢給付金のうち保証期間について支給する給付の現価相当額}を上回
らないものでなければならない。
\end{truefalse}

\begin{answerbox}
\textbf{誤}

  老齢給付金の現価相当額
\end{answerbox}

\begin{truefalse}[年金1 H28 問題1(5)\textcircled{5}]
    \textbf{下線部の内容の正誤を判定し、誤っている場合は正しい内容に改めなさい。}
  \\  

受託保証型確定給付企業年金を除く確定給付企業年金では、一時金として支給する老齢給付
金の額の上限となる現価相当額の計算の基礎となる予定利率は、「一時金を選択したとき」、「老
齢給付金支給開始要件を満たしたとき」、「加入者の資格を喪失したとき」の各時点における\underline{下限予定利率}のうち、最も低い率とする。
\end{truefalse}

\begin{answerbox}
\textbf{誤}

前回の財政計算の日以降の日における下限予定利率のうち最も低い下限予定利率
\end{answerbox}


\begin{shortanswer}[年金1 2025 問題2(1)ア]
確定給付企業年金法第４１条第２項第２号に係る脱退一時金の額は、老齢給付金の全部を年金
として支給するとした場合の老齢給付金（のうち、保証期間について支給する給付）の現価相
当額を上回らないこととされるが、現価相当額の計算に使用する予定利率について簡潔に入力
しなさい。 （250 字以内）
\end{shortanswer}

\begin{answerbox}
前回の財政計算の計算基準日以降の日における下限予定利率のうち、最も低い下限予定利率。
ただし、老齢給付金の額の算定において、加入者の資格を喪失したときから老齢給付金支給開
始要件を満たすまでの期間の全部又は一部について、下限予定利率を下回る利率（当該期間に
応ずる利子に相当する額を加算しない場合にあっては、零）を用いる場合は、当該下回る利率を
用いる期間ごとの当該下回る利率）とする。  
\end{answerbox}

\begin{explanationbox}
  老齢給付金の選択一時金とは予定利率の設定方法が少し違う。
\end{explanationbox}

\newpage

\subsection{給付の額の算定方法１〜DBの給付設計〜}

% TODO 出題数が多いのでサブサブセクションに分けた方がいいかも

\begin{InyouJobunbox}
  \begin{itemize}
    \item 令第24条第1項
    \item 則第25条、第25条の2、第28条第1項
    \item 解釈第3-1\textcircled{4}〜\textcircled{7}
  \end{itemize}
\end{InyouJobunbox}

\begin{truefalse}[年金1 2024 問題1(3)ア]
      \textbf{下線部の内容の正誤を判定し、誤っている場合は正しい内容に改めなさい。}
  \\

  給付額算定用加入者期間は\underline{いかなる場合も加入者期間を上回ることはできない}。
\end{truefalse}

\begin{answerbox}
\textbf{誤}

  端数処理を行う場合は、「給付額算定用加入者期間」（当該端数処理後）は「加入者
期間」（当該端数処理後）を上回ることは差し支えない。
\end{answerbox}

\begin{explanationbox}
  Q\&A項番12

  ただしそのような場合でも端数処理の方法は合理的に定める必要があり、1年以上乖離するような端数処理は認められない。
\end{explanationbox}

\begin{shortanswer}[年金1 2023 問題2(1)イ]
  リスク分担型企業年金を実施する事業主等がその実施事業所を減少させるときに、当該実施
事業所の減少に伴い資格を喪失する加入者に係る調整率を別に定めることができる場合とし
て、確定給付企業年金法施行規則第２５条の２第２項に定める場合を簡潔に入力しなさい。
（250 字以内）
\end{shortanswer}

\begin{answerbox}
  当該実施事業所の減少に伴い、当該リスク分担型企業年金の積立割合、調整率または超過比率が減少すると見込まれる場合。
\end{answerbox}

\begin{truefalse}[年金1 2021 問題1(3)\textcircled{3}]
      \textbf{下線部の内容の正誤を判定し、誤っている場合は正しい内容に改めなさい。}
  \\

リスク分担型企業年金の調整率の改定について、例えば、ある事業年度の決算または財政計
算において「給付財源／調整前給付額の通常の予測に基づく予想額の現価に相当する額」が
0.95 である場合、当該事業年度の末日または当該財政計算の計算基準日の属する事業年度の
\underline{翌事業年度または翌々事業年度の調整率を0.95 とする必要がある}。
\end{truefalse}

\begin{answerbox}
  \textbf{誤}

  翌事業年度または翌々事業年度の調整率を0.95とすることもできるし翌事業年度または翌々事業年度以降5事業年度の調整率を段階的に改定することもできる。
\end{answerbox}

\begin{truefalse}[年金1 2021 問題1(3)\textcircled{4}]
      \textbf{下線部の内容の正誤を判定し、誤っている場合は正しい内容に改めなさい。}
  \\

  キャッシュバランス制度において、給付の額の再評価に用いられる再評価率については、各
期間に用いる再評価率が明確になるように規約に規定する必要があるが、当該率は\underline{上限または
下限を設けることや、規約で定める期間ごとに異}\\
\underline{なる指標を規約に定めて用いること、または
当該率の改定時期を財政再計算の時期とすることは全てできない}。
\end{truefalse}

\begin{answerbox}
  \textbf{誤}

上限または下限を設けることや、規約で定める期間ごとに異なる指標に定めて用いることは可能だが、当該率の改定時期を財政再計算の時期とすることはできない。
\end{answerbox}


\begin{shortanswer}[年金1 2023 問題2(1)エ]
  確定給付企業年金法第３２条第２項において、給付の額は、加入者期間又は当該加入者期間
における給与の額その他これに類するものに照らし、適正かつ合理的な方法により算定され
たものでなければならないと定められている。
「その他これに類するもの」は、ポイント制を採用している場合における労働協約等に定め
られたポイントをいうが、当該ポイントが満たすべき要件として、 『確定給付企業年金法並び
にこれに基づく政令及び省令について（法令解釈）』に掲げる要件を３つ簡潔に入力しなさ
い。 （250 字以内）
\end{shortanswer}

\begin{answerbox}
  \begin{itemize}
    \item 昇格の規定が労働協約等において明確に定められていること
    \item 同一の加入者期間を有する加入者について、最大ポイントと最小ポイントに対する割合に過大な格差がないこと
    \item ポイントは恣意的に決められるものではなく、数理計算が可能であること
  \end{itemize}
\end{answerbox}

\begin{fillblank}[年金1 2020 問題1(1)D]
  \noindent \textbf{◯DB則}
  \begin{lawstructure}
    \begin{lawarticle}{第二十五条の二}{調整率}
      調整率は、リスク分担型企業年金を開始する日の属する事業年度以降の事業年度について、次のとおり定められるものとする。

      （略）

      \lawparagraph{2}
      リスク分担型企業年金を実施する事業主等が、その実施事業所を減少させる場合であって当該減少に伴い当該リスク分担型企業年金の積立割合（調整前給付現価相当額に対する給付財源の割合をいう。以下同じ。）、調整率又は超過比率（調整前給付現価相当額に対する給付財源から調整前給付現価相当額と財政悪化リスク相当額の二分の一の額とを合算した額を控除した額の比率をいう。以下同じ。）が減少すると見込まれるときには、前項の規定にかかわらず、積立割合、調整率又は超過比率が減少しないよう、当該実施事業所の減少に伴い（　　A　　）に係る調整率を別に定めることができる。
    \end{lawarticle}
  \end{lawstructure}
\end{fillblank}

\begin{answerbox}
  \begin{enumerate}[A:]
    \item 資格を喪失する加入者
  \end{enumerate}
\end{answerbox}

\begin{fillblank}[年金2 2022 問題1(3)]
  \noindent \textbf{◯DB則}
  \begin{lawstructure}
    \begin{lawarticle}{第二十五条の二}{調整率}
      調整率は、リスク分担型企業年金を開始する日の属する事業年度以降の事業年度について、次のとおり定められるものとする。

      \lawitem{一}
      リスク分担型企業年金を開始するとき又はリスク分担型企業年金を実施している場合であって給付の設計を変更するとき（掛金の額に係る規約の変更を行う場合に限る。）における調整率は（　　A　　）とする。

      \lawitem{二}
      毎事業年度の決算及び（　　B　　）を行うときに、次に掲げる場合の区分に応じ、次に定める基準を満たすように改定するものとする。

      \hangedlawitem{イ}{3}
      積立金の額に第四十五条第四項に規定するリスク分担型企業年金掛金額の予想額の現価に相当する額を加えた額（以下この条において「給付財源」という。）が調整前給付額の通常の予測に基づく予想額の現価に相当する額（以下この条において「調整前給付現価相当額」という。）に財政悪化リスク相当額（第四十三条第一項に規定する財政悪化リスク相当額をいう。以下この条において同じ。）を加えた額を上回る場合　給付財源と通常予測給付額の現価に相当する額に財政悪化リスク相当額を加えた額が（　　C　　）となること。

      \hangedlawitem{ロ}{3}
      給付財源が調整前給付額の通常の予測に基づく予想額の現価に相当する額を下回る場合　給付財源と通常予測給付額の現価に相当する額が（　　C　　）となること。

      \hangedlawitem{ハ}{3}
      イ及びロ以外の場合　調整率が（　　A　　）となること。

      \lawitem{三}
      前号の調整率の改定は、当該事業年度の末日又は当該財政計算の計算基準日の属する事業年度の翌事業年度又は翌々事業年度以降の事業年度の調整率について行うものとし、当該翌事業年度又は翌々事業年度以降（　　D　　）については、調整率を段階的に引き上げ又は引き下げることができる。

      \lawparagraph{2}
      リスク分担型企業年金を実施する事業主等が、その実施事業所を減少させる場合であって当該減少に伴い当該リスク分担型企業年金の積立割合（調整前給付現価相当額に対する給付財源の割合をいう。以下同じ。）、調整率又は超過比率（調整前給付現価相当額に対する給付財源から調整前給付現価相当額と財政悪化リスク相当額の二分の一の額とを合算した額を控除した額の比率をいう。以下同じ。）が減少すると見込まれるときには、前項の規定にかかわらず、積立割合、調整率又は超過比率が減少しないよう、当該実施事業所の減少に伴い資格を喪失する加入者に係る調整率を（　　E　　）ができる。
    \end{lawarticle}
  \end{lawstructure}
\end{fillblank}

\begin{answerbox}
  \begin{enumerate}[A:]
    \item 一・〇
    \item 財政計算
    \item 同額
    \item 五事業年度
    \item 別に定めること
  \end{enumerate}
\end{answerbox}

\begin{fillblank}[数理人試験 H30 問題4 設問1]
  \noindent \textbf{◯DB則}
  \begin{lawstructure}
    \begin{lawarticle}{第二十五条の二}{調整率}
      調整率は、リスク分担型企業年金を開始する日の属する事業年度以降の事業年度について、次のとおり定められるものとする。

      \lawitem{一}
      リスク分担型企業年金を開始するとき又はリスク分担型企業年金を実施している場合であって給付の設計を変更するとき（掛金の額に係る規約の変更を行う場合に限る。）における調整率は一・〇とする。

      \lawitem{二}
      毎事業年度の決算及び財政計算を行うときに、次に掲げる場合の区分に応じ、次に定める基準を満たすように改定するものとする。

      \hangedlawitem{イ}{3}
      積立金の額に第四十五条第四項に規定するリスク分担型企業年金掛金額の予想額の現価に相当する額を加えた額（以下この条において「給付財源」という。）が調整前給付額の通常の予測に基づく予想額の現価に相当する額（以下この条において「調整前給付現価相当額」という。）に財政悪化リスク相当額（第四十三条第一項に規定する財政悪化リスク相当額をいう。以下この条において同じ。）を加えた額を上回る場合　（　　A　　）となること。

      \hangedlawitem{ロ}{3}
      給付財源が調整前給付額の通常の予測に基づく予想額の現価に相当する額を下回る場合　（　　B　　）となること。

      \hangedlawitem{ハ}{3}
      イ及びロ以外の場合　（　　C　　）となること。

      （以下略）
    \end{lawarticle}
  \end{lawstructure}
\end{fillblank}

\begin{answerbox}
  \begin{enumerate}[A:]
    \item 給付財源と通常予測給付額の現価に財政悪化リスク相当額を加えた額が同額となること
    \item 給付財源と通常予測給付額の現価が同額となること
    \item 調整率が一・〇となること
  \end{enumerate}
\end{answerbox}

\begin{shortanswer}[自作 2026/10/12]
  定給付企業年金法施行規則第２５条第１号において、給付の額を算定する方法として定額制（令２４条第１項第１号）、給与比例制（令２４条第１項第２号）、CB制（令２４条第１項第３号）による算定結果を「組み合わせた方法」を採用することができると定められている。「組み合わせた方法」として『確定給付企業年金法並びにこれに基づく政令及び省令について（法令解釈）』に記載されている4つの方法を簡記しなさい。
\end{shortanswer}

\begin{answerbox}
  \begin{itemize}
    \item 加法
    \item 減法（ただし減法の結果、給付の額が零を上回るものとする）
    \item 一定の数値を乗じる方法または一定の数値で除する方法
    \item 加入者期間、給付額算定用加入者期間、資格喪失事由、労働協約等に定める職種等または年齢に応じて異なる算定方法とする方法
  \end{itemize}
\end{answerbox}

\begin{truefalse}[年金2 2021 問題1(1)\textcircled{1}]
      \textbf{「確定給付企業年金に関する数理実務ガイダンス」に記載されているリスク分担型企業年
金の取り扱いとして適当なものを選択しなさい。}
  \\  

  \begin{enumerate}
    \item 財政方式によらず、通常予測給付現価の算定対象に将来加入者を含むことを基本とする。
    \item 非継続基準は結果的に1.00 となるが、最低積立基準額は計算する必要がある。
    \item 積立上限額の計算を要しない。
    \item 承継事業所償却積立金は設定できない。
  \end{enumerate}
\end{truefalse}

\begin{answerbox}
  \textbf{1, 2, 3, 4}
\end{answerbox}

\begin{truefalse}[年金2 2021 問題1(1)\textcircled{2}]
      \textbf{下線部の内容の正誤を判定し、誤っている場合は正しい内容に改めなさい。}
  \\  

  「確定給付企業年金に関する数理実務ガイダンス」では、給付の額の算定において、給付の
再評価又は額の改定を行う場合にあっては、当該再評価及び額の改定に用いる指標の予測を、
当該指標に係る実績及び将来の見通しに基づいて\underline{事業主等が主体的に}決定することと規定さ
れている。
\end{truefalse}

\begin{answerbox}
  \textbf{正}
\end{answerbox}

\newpage

\subsection{給付の額の算定方法２〜規約で定める数値〜}

\begin{InyouJobunbox}
  \begin{itemize}
    \item 令第24条第2項
    \item 則第26条〜第27条
    \item 解釈第3-1\textcircled{6}\textcircled{8}\textcircled{9}
  \end{itemize}
\end{InyouJobunbox}

\begin{fillblank}[年金1 2019 問題1(1)DEF]
  \noindent \textbf{◯DB令}
  \begin{lawstructure}
    \begin{lawarticle}{第二十四条}{給付の額の算定方法}
      法第三十二条第二項の政令で定める方法は、次の各号のいずれかに該当する方法とする。

      （略）

      \lawparagraph{2}
      前項第一号から第三号までに規定する規約で定める数値は、厚生労働省令で定めるところにより、支給開始時における（　　A　　）、支給期間、保証期間（保証期間を定めた場合に限る。）その他厚生労働省令で定めるものに応じたものとしなければならない。

      （略）
    \end{lawarticle}

    \noindent \textbf{◯DB則}
    \begin{lawarticle}{第二十七条}{規約で定める数値のその他の算定基礎}
      令第二十四条第二項の厚生労働省令で定めるものは、次のとおりとする。

      （略）

      \lawitem{一}
      加入者の資格を喪失した者が当該資格を喪失したときから（　　B　　）を満たすまでの期間（老齢給付金の額に当該期間に応ずる利子に相当する額を加算することとなっている場合に限る。この場合において、当該利子については前条第三項第一号の規定を適用しない。）

      \lawitem{二}
      老齢給付金の受給権者が死亡した場合にその遺族（法第四十八条に規定する遺族給付金（法第二十九条第二項第二号に規定する遺族給付金をいう。以下同じ。）を受けることができる遺族をいう。以下同じ。）に支給される遺族給付金の給付の設計（老齢給付金の受給権の裁定のときに、当該老齢給付金の受給権者の死亡によりその遺族に支給されるべき遺族給付金の給付の設計を選択できる場合に限る。）

      \lawitem{三}
      加入者の資格を喪失した事由

      \lawitem{四}
      加入者の資格を喪失した日における当該加入者の年齢

      \lawitem{五}
      （　　C　　）

    \end{lawarticle}
  \end{lawstructure}
\end{fillblank}

\begin{answerbox}
  \begin{enumerate}[A:]
    \item 受給権者の年齢
    \item 老齢給付金受給開始要件
    \item 加入者である期間（または加入者期間）
  \end{enumerate}
\end{answerbox}

\newpage

\subsection{給付の額の算定方法３〜額の改定〜}

\begin{InyouJobunbox}
  \begin{itemize}
    \item 令第24条第3項
    \item 則第28条第2項
    \item 解釈第3-1\textcircled{6}\textcircled{10}
  \end{itemize}
\end{InyouJobunbox}

\begin{shortanswer}[年金1 2020 問題2(1)\textcircled{2}]
   確定給付企業年金法施行令第 24 条第 3 項に定めるとおり、年金として支給する給付の額は、
当該給付が支給される間において、規約に定めることにより改定することができるが、確定給
付企業年金法施行規則第28 条第2 項に定める「額を改定する方法」を３つ挙げなさい。
\end{shortanswer}

\begin{answerbox}
  （以下は一例）
  \begin{itemize}
    \item 給付の支給を開始して一定の期間が経過したときまたは一定の年齢に達したときに、定率を乗じることで改定する方法
\item 規約で定める期間ごとに、前の期間の給付の額に、当該前の期間の給付の額に指標を乗じて得た額を加算することで改定する方法
\item 給付の支給を開始した後に加入者期間の全部または一部により給付の額を改定する方法
  \end{itemize}
\end{answerbox}

\begin{fillblank}[数理人試験 2025 問題4 設問1]
  \noindent \textbf{◯DB則}
  \begin{lawstructure}
    \begin{lawarticle}{第二十八条}{給付の額の再評価等の方法}
      令第二十四条第一項第三号の再評価は、規約で定める期間ごとに、次条第一項各号に掲げるもの（以下「指標」という。）を用いて行うものとする。
      
      \lawitem{一}
      給付の支給を開始して一定の期間が経過したとき又は一定の年齢に達したときに、次のいずれかの方法により改定する方法

      \hangedlawitem{イ}{3}
      （　　A　　）

      \hangedlawitem{ロ}{3}
      令第二十四条第一項各号のいずれかの方法（当該給付の額を算定した方法を除く。）

      \lawitem{二}
      規約で定める期間ごとに、次のいずれかの加算を行うことにより改定する方法

      \hangedlawitem{イ}{3}
      前の期間の給付の額に、当該前の期間の給付の額に（　　B　　）を加算すること。

      \hangedlawitem{ロ}{3}
      （　　C　　）に、規約で定める期間、指標を第二十六条第三項第一号の予定利率とみなして算定するとした場合における給付の額が（　　C　　）を上回る額その他これに類する額を加算すること（当該指標が第二十六条第三項第一号の予定利率を上回る場合に限る。）。

      \lawitem{三}
      給付の支給を開始した後に（　　D　　）により給付の額を改定する方法

    \end{lawarticle}
  \end{lawstructure}  
\end{fillblank}

\begin{answerbox}
  \begin{enumerate}[A:]
    \item 定率を乗じる方法
    \item 指標を乗じて得た額
    \item あらかじめ定めた給付の額
    \item 加入者期間の全部又は一部
  \end{enumerate}
\end{answerbox}

\begin{explanationbox}
  Dの類題：年金1 2025 問題1(3)ウ
\end{explanationbox}

\begin{fillblank}[数理人試験 H28 問題3 設問1]
  \noindent \textbf{◯DB則}
  \begin{lawstructure}
    \begin{lawarticle}{第二十八条}{給付の額の再評価等の方法}
      令第二十四条第一項第三号の再評価は、規約で定める期間ごとに、次条第一項各号に掲げるもの（以下「指標」という。）を用いて行うものとする。
      
      \lawitem{一}
      給付の支給を開始して一定の期間が経過したとき又は一定の年齢に達したときに、次のいずれかの方法により改定する方法

      \hangedlawitem{イ}{3}
      （　　A　　）

      \hangedlawitem{ロ}{3}
      令第二十四条第一項各号のいずれかの方法（当該給付の額を算定した方法を除く。）

      \lawitem{二}
      規約で定める期間ごとに、次のいずれかの加算を行うことにより改定する方法

      \hangedlawitem{イ}{3}
      （　　B　　）に、当該（　　B　　）に（　　C　　）を加算すること。

      \hangedlawitem{ロ}{3}
      （　　D　　）に、規約で定める期間、指標を第二十六条第三項第一号の予定利率とみなして算定するとした場合における給付の額が（　　D　　）を上回る額その他これに類する額を加算すること（当該指標が第二十六条第三項第一号の予定利率を上回る場合に限る。）。

      \lawitem{三}
      給付の支給を開始した後に加入者期間の全部又は一部により給付の額を改定する方法

    \end{lawarticle}
  \end{lawstructure}  
\end{fillblank}

\begin{answerbox}
  \begin{enumerate}[A:]
    \item 定率を乗じる方法
    \item 前の期間の給付の額
    \item 指標を乗じて得た額
    \item あらかじめ定めた給付の額
  \end{enumerate}
\end{answerbox}

\begin{shortanswer}[数理人試験 H28 問題3 設問2]
  給付の額の改定方法として、確定給付企業年金法施行規則第二十八条第二項第二号ロに規定す
る方法に関し、事務連絡「＜参考１＞規約型確定給付企業年金規約例」において、２通りの方
法が記載されている。当該２通りの方法の内容を簡記せよ。
\end{shortanswer}

\begin{answerbox}
  ＜指標に応じた年金現価率を用いる方法＞

改定後の給付の額はあらかじめ規約に定めた給付の額に、老齢給付金の支給要件を満たした日の仮想個人勘定残高を改定時に適用される指標に応じて規約に定める年金現価率で除して得た額があらかじめ規約に定めた給付の額を上回る額を加算した額とする。

＜指標に応じた残余支給期間に係る年金現価率を用いる方法＞

改定後の給付の額は、あらかじめ規約に定めた給付の額に、改定前の給付の額の算定に用いた指標および残余支給期間に応じて規約に定める年金現価率を乗じて得た額を改定時に適用される指標および残余支給期間に応じて規約に定める年金現価率で除して得た額があらかじめ規約に定めた給付の額を上回る額を加算した額とする。
\end{answerbox}

\begin{shortanswer}[数理人試験 H28 問題3 設問3]
  下記前提を用いて設問２の２通りの方法によるX＋15年度の年金額をそれぞれ算出せよ。解答
にあたっては、金額の端数処理は円未満を四捨五入することとし、設問２において解答した順
に解答せよ。なお、解答に至るまでの計算式や過程は論述しなくてよい。

＜前提＞
\begin{itemize}
  \item 支給開始時（X年度期初）の仮想個人勘定残高：1,000万円
  \item 20年確定年金（年1回期末払い）
  \item 受給中の死亡時には遺族に転給を行うものとする
  \item 給付の額の改定は年1回期初に行うものとする
  \item 設問２の方法以外の給付の額の改定は行わないものとする
  \item 遺族一時金・選択一時金を選択せず、年金の給付を選択するものとする
  \item 指標の実績値は以下のとおり
  
  \begin{center}
    X年度～X＋14年度：2.0％、X＋15年度～X＋19年度：2.5％
  \end{center}

  \item 年１回期末払い確定年金現価率は以下のとおり
  
  \begin{table}[htbp]
    \centering
    \begin{tabular}{|l|c|c|} \hline
      & 2.0％ & 2.5％ \\ \hline
      ５年確定年金現価率 & 4.71346 & 4.64583 \\ \hline
      １５年確定年金現価率 & 12.84926 & 12.38138 \\ \hline 
      ２０年確定年金現価率 & 16.35143 & 15.58916 \\ \hline
    \end{tabular}
  \end{table}
\end{itemize}
\end{shortanswer}

\begin{answerbox}
＜指標に応じた年金現価率を用いる方法＞

（改定前年金額）1000万円÷16.35143=611,567円

（比較用の年金額）1000万円÷15.58916=641,471円

（改定後年金額）611,567+Max(0, 641,471-611,567)=641,741円

＜指標に応じた残余支給期間に係る年金現価率を用いる方法＞

（改定前年金額）1000万円÷16.35143=611,567円

（改定後年金額）611,567×(4.71346÷4.63583)=620,470円
\end{answerbox}

\newpage 

\subsection{給付の額の算定方法４〜指標〜}

\begin{InyouJobunbox}
  \begin{itemize}
    \item 令第24条第4項
    \item 則第29条、第80条
    \item 解釈第3-3
  \end{itemize}
\end{InyouJobunbox}

\begin{shortanswer}[年金1 2022 問題2(1)ウ]
  確定給付企業年金法施行規則第29 条において、キャッシュバランス制度の再評価率について
定められているが、 「その他の客観的な指標であって、合理的に予測することが可能なもの」
として、『確定給付企業年金法並びにこれに基づく政令及び省令について（法令解釈）』に例
示されているものを3 つ簡潔に入力しなさい。
\end{shortanswer}

\begin{answerbox}
  \begin{itemize}
    \item 総務省において作成する年平均の全国消費者物価指数
    \item 厚生労働省において作成する年平均の賃金指数
    \item 確定給付企業年金法施行規則第80条第3項第1号または第2号に掲げる有価証券指標（東証株価指数、Russell/Nomura Prime インデックス）
  \end{itemize}
  
\end{answerbox}

\begin{truefalse}[年金1 2018 問題1(1)\textcircled{1}]
    \textbf{下線部の内容の正誤を判定し、誤っている場合は正しい内容に改めなさい。}
  \\

  確定給付企業年金法施行規則第 29 条において、令第24 条第 4 項に規定する給付の額の再
評価等に用いる率は次のとおり定められている。
\begin{enumerate}
  \item 定率
\item 国債の利回りその他の\underline{客観的な指標}
\item  積立金の運用利回りの実績
\item  前３号に掲げる率を組み合わせたもの
\item  前３号に掲げる率にその上限又は下限を定めたもの
\end{enumerate}
\end{truefalse}

\begin{answerbox}
  \textbf{誤}

  客観的な指標であって、合理的に予測することが可能なもの
\end{answerbox}

\begin{truefalse}[年金2 2021 問題1(1)\textcircled{2}]
      \textbf{下線部の内容の正誤を判定し、誤っている場合は正しい内容に改めなさい。}
  \\

  「確定給付企業年金に関する数理実務ガイダンス」では、給付の額の算定において、給付の
再評価又は額の改定を行う場合にあっては、当該再評価及び額の改定に用いる指標の予測を、
当該指標に係る実績及び将来の見通しに基づいて\underline{事業主等が主体的に}決定することと規定さ
れている。
\end{truefalse}

\begin{answerbox}
  \textbf{正}
\end{answerbox}

\begin{fillblank}[数理人試験 H29 問題1 設問2]
  \noindent \textbf{◯DB則}
  \begin{lawstructure}
    \begin{lawarticle}{第二十九条}{給付の額の再評価等に用いる率}
      令第二十四条第四項に規定する厚生労働省令で定めるものは、次のとおりとする。
      ただし、同条第一項第三号に掲げる給付の額の算定方法を用いて給付の額を計算する場合にあっては、
      次の各号のいずれの率に基づき再評価を行う場合でも、当該再評価後の（　　A　　）が、
      当該再評価を行わなかった場合の（　　A　　）を下回ってはならない。

      \lawitem{一}
      （　　B　　）

      \lawitem{二}
      （　　C　　）の利回りその他の客観的な指標であって、合理的に予測することが可能なもの

      \lawitem{三}
      積立金の運用利回りの（　　D　　）

      \lawitem{四}
      前三号に掲げる率を組み合わせたもの

      \lawitem{五}
      前三号に掲げる率にその上限又は下限を定めたもの
      
    \end{lawarticle}

    \noindent \textbf{◯法令解釈}

    \begin{lawarticle}{}{}
      
      第３ 給付の額に関する事項

      \lawparagraph{1}
      （略）

      \lawparagraph{2}
      （略）

      \lawparagraph{3}
      規則第２９条第１項第２号に「その他の客観的な指標であって、合理的に予測する
ことが可能なもの」とあるが、例えば、次に掲げるものはこれに該当するものである
こと。

      \lawitem{\textcircled{1}}
      総務省において作成する年平均の（　　E　　）

      \lawitem{\textcircled{2}}
      厚生労働省において作成する年平均の（　　F　　）

      \lawitem{\textcircled{3}}
      規則第８０条第３項第１号又は第２号に掲げる有価証券指標

    \end{lawarticle}

  \end{lawstructure}  
\end{fillblank}

\begin{answerbox}
  \begin{enumerate}[A:]
    \item 累計額
    \item 定率
    \item 国債
    \item 実績
    \item 賃金指数
    \item 全国消費者物価指数
  \end{enumerate}
\end{answerbox}

\newpage

\subsection{給付の通則}

\begin{InyouJobunbox}
  \begin{itemize}
    \item 法第33条〜第35条
    \item 令第25条〜第26条
    \item 則第34条
  \end{itemize}
\end{InyouJobunbox}

\begin{truefalse}[年金1 H28 問題2(1)\textcircled{3}]
    \textbf{下線部の内容の正誤を判定し、誤っている場合は正しい内容に改めなさい。}
  \\

   「死亡した受給権者が死亡前に請求していなかった未請求の給付」または「死亡した受給権
者に支給すべき未支給給付」がある場合、死亡した受給権者の遺族のうち規約で定めるものは、
自己の名でその給付を請求することができる。
\end{truefalse}

\begin{answerbox}
  \textbf{正}
\end{answerbox}

\begin{fillblank}[自作 2026/10/12]
  \noindent \textbf{◯DB法}
  \begin{lawstructure}
    \begin{lawarticle}{第三十三条}{年金給付の支給期間等}
      年金給付の支給期間及び支払期月は、政令で定める基準に従い規約で定めるところによる。ただし、終身又は（　　A　　）年以上にわたり、毎年一回以上定期的に支給するものでなければならない。
    \end{lawarticle}
  \end{lawstructure}

  \noindent \textbf{◯DB令}
  \begin{lawstructure}
    \begin{lawarticle}{第二十五条}{支給期間及び支払期月}
      第三十三条の政令で定める基準は、次のとおりとする。
\lawitem{一} 保証期間を定める場合にあっては、（　　B　　）年を超えない範囲内で定めること。
\lawitem{二} 年金給付の支払期月は、毎年一定の時期であること。

    \end{lawarticle}
  \end{lawstructure}
\end{fillblank}

\begin{answerbox}
  \begin{enumerate}[A:]
    \item 5
    \item 20
  \end{enumerate}
\end{answerbox}

\begin{truefalse}[年金1 2024 問題1(3)ウ]
      \textbf{下線部の内容の正誤を判定し、誤っている場合は正しい内容に改めなさい。}
  \\

  受給権は、譲り渡し、担保に供し、又は差し押さえることができない。ただし、\underline{老齢給付金お
よび障害給付金}を受ける権利を国税滞納処分（その例による処分を含む。）により差し押さえる
場合は、この限りでない。
\end{truefalse}

\begin{answerbox}
  \textbf{誤}

  老齢給付金、脱退一時金及び遺族給付金
\end{answerbox}

\newpage

\subsection{老齢給付金１}


\begin{InyouJobunbox}
  \begin{itemize}
    \item 法第36条〜第37条
    \item 令第28条
    \item 解釈第3-1\textcircled{11}
  \end{itemize}
\end{InyouJobunbox}

\begin{fillblank}[年金1 2020 問題1(1)A、年金1 2019 問題1(1)BC]
  \noindent \textbf{◯DB法}
  \begin{lawstructure}
    \begin{lawarticle}{第三十六条}{支給要件}
      老齢給付金は、加入者又は加入者であった者が、規約で定める老齢給付金を受けるための要件を満たすこととなったときに、その者に支給するものとする。

      \lawparagraph{2}
      前項に規定する規約で定める要件は、次に掲げる要件（第四十一条第二項第二号において「老齢給付金支給開始要件」という。）を満たすものでなければならない。

      \lawitem{一}
      （　　A　　）の規約で定める年齢に達したときに支給するものであること。
      
      \lawitem{二}
      政令で定める年齢以上前号の規約で定める年齢未満の規約で定める年齢に達した日以後に実施事業所に使用されなくなったときに支給するものであること（規約において当該状態に至ったときに老齢給付金を支給する旨が定められている場合に限る。）。
    
      \lawparagraph{3}
      前項第二号の政令で定める年齢は、（　　B　　）未満であってはならない。

      \lawparagraph{4}
      規約において、二十年を超える加入者期間を老齢給付金の給付を受けるための要件として定めてはならない。
    \end{lawarticle}
  \end{lawstructure}  
\end{fillblank}

\begin{answerbox}
  \begin{enumerate}[A:]
    \item 六十歳以上七十歳以下
    \item 五十歳
  \end{enumerate}
\end{answerbox}

\begin{truefalse}[年金1 2022 問題1(3)ア]
      \textbf{下線部の内容の正誤を判定し、誤っている場合は正しい内容に改めなさい。}
  \\

  老齢給付金の受給権者が支給の繰下げを申し出たときに使用する繰下利率は、\underline{資格喪失事
由、資格喪失時の年齢、}\\
\underline{加入者であった期間、加入者であった期間におけ
る給与の額、申し
出から年金の支給が開始されるまでの期間におけ}\\
\underline{る年齢等}に基づき合理的に定める必要があ
る。
\end{truefalse}

\begin{answerbox}
  \textbf{誤}

  資格喪失事由、資格喪失時の年齢、労働協約等に定める職種等、加入者であった期間、申し出から年金の支給が開始されるまでの期間における年齢等
\end{answerbox}

\begin{truefalse}[年金1 H29 問題2(1)\textcircled{1}]
        \textbf{以下の文章の正誤を判定し、誤っている場合はその理由を記述しなさい。}
  \\

  規約で定めた場合には、老齢給付金の支給を受ける年齢に達した後に加入者期間に係る要件
を満たしたときにも老齢給付金を支給することができる。
\end{truefalse}

\begin{answerbox}
  \textbf{正}
\end{answerbox}

\begin{truefalse}[年金1 H28 問題2(1)\textcircled{4}]
\textbf{以下の文章の正誤を判定し、誤っている場合はその理由を記述しなさい。}
  \\

   「加入者期間 2 年以上の者が 60 歳の誕生日の属する月の末日を迎えたとき」を老齢給付金の
支給要件とすることができる。
\end{truefalse}

\begin{answerbox}
  \textbf{誤}

  「誕生日の属する月の末日」等の年齢以外の要件を定めることはできない。
\end{answerbox}

\begin{explanationbox}
  Q\&A項番25、規約令規第21条留意事項

  年齢要件は「年数」のみ指定可能。「誕生日の属する月の末日」のような指定のほか、「64歳と6月」のように月数を指定することも不可。
\end{explanationbox}

\begin{truefalse}[年金1 2025 問題1(3)エ\textcircled{1}]
\textbf{以下の文章の正誤を判定し、誤っている場合はその理由を記述しなさい。}
  \\

  老齢給付金の支給要件について、「加入者期間２０年以上の者が６０歳以降最初に到来する
３月３１日を迎えたとき」とするような給付設計は設計可能である。
\end{truefalse}

\begin{answerbox}
  \textbf{誤}

  「誕生日の属する月の末日」等の年齢以外の要件を定めることはできない。
\end{answerbox}

\begin{truefalse}[年金1 2025 問題1(3)エ\textcircled{2}]
\textbf{以下の文章の正誤を判定し、誤っている場合はその理由を記述しなさい。}
  \\

老齢給付金の支給を受ける年齢に達した後に加入者期間に係る要件を満たしたときにも老
齢給付金を支給すること
\end{truefalse}

\begin{answerbox}
  \textbf{正}
\end{answerbox}


\newpage

\subsection{老齢給付金２}


\begin{InyouJobunbox}
  \begin{itemize}
    \item 法第38条〜第40条
    \item 令第29条〜第30条
    \item 則第30条、第35条
    \item 解釈第3-4
  \end{itemize}
\end{InyouJobunbox}


\begin{fillblank}[年金1 H28 問題1(1)]
  \noindent \textbf{◯DB法}
  \begin{lawstructure}
    \begin{lawarticle}{第三十八条}{支給の方法}
      老齢給付金は、年金として支給する。

      \lawparagraph{2}
      老齢給付金は、規約でその（　　A　　）を一時金として支給することができることを定めた場合には、前項の規定にかかわらず、政令で定める基準に従い規約で定めるところにより、一時金として支給することができる。
    \end{lawarticle}
  \end{lawstructure}
  
  \noindent \textbf{◯DB令}
  \begin{lawstructure}
    \begin{lawarticle}{第二十九条}{老齢給付金を一時金として支給する場合の基準}
      法第三十八条第二項の政令で定める基準は、次のとおりとする。

      \lawitem{一}
        年金として支給する老齢給付金について（　　B　　）が定められていること。
      
      \lawitem{二}
      老齢給付金の受給権者の選択により一時金として支給するものであること。

      \lawitem{三}
      前号の選択は、法第三十条第一項の（　　C　　）に併せて行うとき、又は年金として支給する老齢給付金の支給を開始してから五年を経過した日以後に行うときに限り、することができるものであること。ただし、年金として支給する老齢給付金の受給権者に厚生労働省令で定める特別の事情がある場合にあっては、当該老齢給付金の支給を開始してから五年を経過する日までの間においても、同号の選択をすることができるものであること。

    \end{lawarticle}
  \end{lawstructure}

  \noindent \textbf{◯DB則}
  \begin{lawstructure}
    \begin{lawarticle}{第三十条}{老齢給付金について一時金を選択することができる特別の事情}
      令第二十九条第三号の厚生労働省令で定める特別の事情は、次のとおりとする。

      \lawitem{一}
      受給権者又はその属する世帯の（　　D　　）する者が、震災、風水害、火災その他これらに類する災害により、住宅、家財又はその他の財産について著しい損害を受けたこと。

      \lawitem{二}
      受給権者がその（　　E　　）することが困難であること。

      \lawitem{三}
      受給権者が心身に重大な障害を受け、又は長期間入院したこと。

      \lawitem{四}
      その他前三号に準ずる事情

    \end{lawarticle}
  \end{lawstructure}
\end{fillblank}

\begin{answerbox}
  \begin{enumerate}[A:]
    \item 全部又は一部
    \item 保証期間
    \item 請求
    \item 生計を主として維持
    \item 債務を弁済
  \end{enumerate}
\end{answerbox}

\begin{fillblank}[自作 2026/10/12]
  \noindent \textbf{◯DB法}
  \begin{lawstructure}
    \begin{lawarticle}{第三十九条}{支給停止}
      老齢給付金の受給権者が、（　　A　　）を支給されたときは、第三十六条第一項の規定にかかわらず、政令で定める基準に従い規約で定めるところにより、老齢給付金の額の全部又は一部につき、その支給を停止することができる。
    \end{lawarticle}
  \end{lawstructure}

  \noindent \textbf{◯DB令}
  \begin{lawstructure}
    \begin{lawarticle}{第三十条}{老齢給付金の支給停止の基準}
      法第三十九条の政令で定める基準は、次のとおりとする。

      \lawitem{一}
      まだ支給されていない老齢給付金の現価相当額が（　　A　　）の現価相当額を超える場合における当該超える部分については、支給を停止しないこと。
      
      \lawitem{二}
（　　A　　）の支給期間が終了したときに老齢給付金の支給期間が終了していない場合には、当該（　　A　　）の支給期間が終了した後の老齢給付金の支給期間については、支給を停止しないこと。

\lawparagraph{2}
  第二十三条第四項の規定は、前項第一号の現価相当額を計算する場合について準用する。


    \end{lawarticle}
  \end{lawstructure}
\end{fillblank}

\begin{answerbox}
  \begin{enumerate}[A:]
    \item 障害給付金
  \end{enumerate}
\end{answerbox}

\begin{shortanswer}[年金1 2024 問題2(1)イ]
  年金として支給する老齢給付金の支給を開始してから５年を経過する日までの間において、年
金として支給する老齢給付金の受給権者に厚生労働省で定める特別の事情がある場合に老齢
給付金を一時金として支給できるが、この特別の事情について確定給付企業年金法施行規則第
３０条に定められているものを４つ簡潔に入力しなさい。 （250 字以内）
\end{shortanswer}

\begin{answerbox}
  \begin{itemize}
    \item 受給権者又はその属する世帯の生計を主として維持する者が、震災、風水害、火災その他これら
に類する災害により、住宅、家財又はその他の財産について著しい損害を受けたこと。
\item 受給権者がその債務を弁済することが困難であること。
\item 受給権者が心身に重大な障害を受け、又は長期間入院したこと。
\item その他前三号に準ずる事情
  \end{itemize}
\end{answerbox}

\newpage

\subsection{脱退一時金}

\begin{InyouJobunbox}
  \begin{itemize}
    \item 法第41条、第42条
    \item 令第27条
    \item 解釈第3-1\textcircled{11}
  \end{itemize}
\end{InyouJobunbox}

\begin{fillblank}[年金1 2020 問題1(1)BC、自作]
  \noindent \textbf{◯DB令}
  \begin{lawstructure}
    \begin{lawarticle}{第二十七条}{脱退一時金の支給要件及び失権}
      法第四十一条第二項第二号に係る脱退一時金を受けるための要件を規約で定める場合にあっては、次に掲げる要件を満たすものでなければならない。

      \lawitem{一}
      当該加入者が老齢給付金の受給権者となったときに支給する老齢給付金の全部又は一部に代えて支給するものであること。

      \lawitem{二}
      当該老齢給付金に（　　A　　）が定められていること。

      \lawitem{三}
      当該加入者の（　　B　　）により当該脱退一時金の全部の支給の繰下げができるものであること。

      \lawparagraph{2}
      法第四十一条第四項の規定に基づき全部又は一部の支給の繰下げの申出をした脱退一時金の受給権は、次の各号のいずれかに該当することとなったときは、消滅する。
      
      \lawitem{一}
      脱退一時金の受給権者が（　　C　　）したとき。

      \lawitem{二}
      脱退一時金の受給権者（法第四十一条第二項第二号に係る脱退一時金の受給権者に限る。）が（　　D　　）の受給権者となったとき。

      \lawitem{三}
      再加入者となる前に当該確定給付企業年金の脱退一時金の受給権者となった者について、当該再加入者の当該確定給付企業年金における前後の加入者期間を合算したとき。
    \end{lawarticle}
  \end{lawstructure}  
\end{fillblank}

\begin{answerbox}
  \begin{enumerate}[A:]
    \item 保証期間
    \item 選択
    \item 死亡
    \item 老齢給付金
  \end{enumerate}
\end{answerbox}

\begin{truefalse}[年金1 2022 問題1(3)エ]
        \textbf{下線部の内容の正誤を判定し、誤っている場合は正しい内容に改めなさい。}
  \\

  脱退一時金を受けるための要件を満たす者（\underline{死亡したとき}を除く）は、規約で定めるところ
により、事業主や企業年金基金に脱退一時金の全部又は一部の支給の繰下げの申し出をする
ことができる。
\end{truefalse}

\begin{answerbox}
  \textbf{誤}

  その使用される事業所または船舶が実施事業所ではなくなったとき
\end{answerbox}

\begin{truefalse}[年金1 2019 問題1(2)\textcircled{1}]
        \textbf{下線部の内容の正誤を判定し、誤っている場合は正しい内容に改めなさい。}
  \\

脱退一時金を受けるための要件を満たす者は、\underline{厚生年金保険の被保険者でなくなったこと}
により加入者の資格を喪失した場合を除き、規約で定めるところにより、事業主等に当該脱
退一時金の全部または一部の支給の繰下げの申出をすることができる。
\end{truefalse}

\begin{answerbox}
  \textbf{誤}

  その使用される事業所または船舶が実施事業所ではなくなったとき
\end{answerbox}

\begin{truefalse}[年金1 H29 問題1(5)\textcircled{5}]
        \textbf{下線部の内容の正誤を判定し、誤っている場合は正しい内容に改めなさい。}
  \\

リスク分担型企業年金において、脱退一時金を繰り下げた場合の脱退一時金の額には、\underline{繰下
げを終了したとき}の調整率を適用する。
\end{truefalse}

\begin{answerbox}
  \textbf{正}
\end{answerbox}

\begin{truefalse}[年金1 H28 問題2(1)\textcircled{2}]
\textbf{以下の文章の正誤を判定し、誤っている場合はその理由を記述しなさい。}
  \\

  規約で定めるところにより、脱退一時金を受けるための要件を満たす全ての者が、脱退一時
金の支給の繰下げを申し出ることができる。
\end{truefalse}

\begin{answerbox}
  \textbf{誤}

  その使用される事業所または船舶が実施事業所で亡くなったことにより脱退一時金を受けるための要件を満たした者は脱退一時金の支給の繰下げを申し出ることができない。
\end{answerbox}

\begin{shortanswer}[年金1 2024 問題2(1)エ]
  確定給付企業年金法第４１条第２項第２号に係る脱退一時金（老齢給付金を受けるための要件
のうち老齢給付金支給開始要件以外の要件を満たすものに支給するもの）を受けるための要件
を規約に定める場合にあって、満たすべき要件として確定給付企業年金法施行令第２７条第１
項に定められているものを３つ簡潔に入力しなさい。 （250 字以内）
\end{shortanswer}

\begin{answerbox}
  \begin{itemize}
    \item 当該加入者が老齢給付金の受給権者となったときに支給する老齢給付金の全部又は一部に代え
て支給するものであること。
\item 当該老齢給付金に保証期間が定められていること。
\item 当該加入者の選択により当該脱退一時金の全部の支給の繰下げができるものであること。
  \end{itemize}
\end{answerbox}

\newpage

\subsection{障害給付金}

\begin{InyouJobunbox}
  \begin{itemize}
    \item 法第43条〜第46条
    \item 令第31条〜第32条
  \end{itemize}
\end{InyouJobunbox}

\begin{truefalse}[年金1 H28 問題2(1)\textcircled{1}]
  \textbf{以下の文章の正誤を判定し、誤っている場合はその理由を記述しなさい。}
  \\

  「年金として支給する障害給付金」「一時金として支給する障害給付金」はともに支給の繰下
げを行うことができない。
\end{truefalse}

\begin{answerbox}
  \textbf{正}
\end{answerbox}



\begin{truefalse}[年金1 2025 問題1(3)エ\textcircled{3}]
\textbf{以下の文章の正誤を判定し、誤っている場合はその理由を記述しなさい。}
  \\

障害給付金について支給の繰下げを可能とするような給付設計は設計可能である。
\end{truefalse}

\begin{answerbox}
  \textbf{誤}

  障害給付金は支給繰下げできない
\end{answerbox}

\newpage

\subsection{遺族給付金}



\begin{InyouJobunbox}
  \begin{itemize}
    \item 法第47条〜第51条
    \item 令第33条
  \end{itemize}
\end{InyouJobunbox}

\begin{truefalse}[年金1 2020 問題1(3)\textcircled{4}]
          \textbf{下線部の内容の正誤を判定し、誤っている場合は正しい内容に改めなさい。}
  \\

  老齢給付金または障害給付金（以下「老齢給付金等」 ）の給付を受けている者が死亡したとき
にその遺族に対し年金として支給する遺族給付金の支給期間については、当該老齢給付金等の
支給期間として\underline{規約において一定の期間を定めていた場合は 3 年未満とすることができる}。た
だし、当該老齢給付金等の支給期間のうち給付を受けていない期間を下回ることができない。
\end{truefalse}

\begin{answerbox}
  \textbf{誤}

  規約において一定の期間を定めていた場合は5年未満とすることができる
\end{answerbox}

\newpage

\subsection{給付の制限}

出題例なし